% SOURCE_AUTHORITY: sga5_fr_workpass.tex, lines 15276--15286 (LF-normalized unit).
% SOURCE_SHA256: 791F4EFFC5E02832D5D77ED03518C8156D6F07E4C8238B03545DB93D883FBB28
Demostración del lema 1: se puede (sustituyendo, si es preciso, $L$ por un complejo cuasi-isomorfo) suponer que el complejo $L$, y por tanto también $T(L)$, es proyectivo en cada grado, lo que supondremos en lo sucesivo; se puede, por consiguiente, suponer que el isomorfismo $u$ de $D^-(A_0)$ procede de un cuasi-isomorfismo $T(L)\to M$, denotado asimismo por $u$. Primero nos reduciremos al caso en que $u$ sea sobreyectivo. El complejo $N_0$ definido por
\[
N_0^i=M^i\oplus M^{i-1},
\quad
 d^i(x,y)=(d^i x,x-d^{i-1}y)
\]
($x\in M^i$, $y\in M^{i-1}$) es manifiestamente acíclico y proyectivo en cada grado; se define un epimorfismo de complejos $N_0\to M$ mediante la proyección $N_0^i\to M^i$ en grado $i$. El corolario del lema 3 permite levantar $N_0$ a un complejo $N$ acíclico, acotado por la derecha y proyectivo en cada grado. Se puede entonces sustituir $L$ por $L\oplus N$, que es cuasi-isomorfo a él (pues $N$ es acíclico) y proyectivo en cada grado; al mismo tiempo se sustituye $u$ por el cuasi-isomorfismo
\[
T(L\oplus N)=T(L)\oplus N_0\longrightarrow M
\]
definido por $u$ sobre $T(L)$ y por el epimorfismo $N_0\to M$ sobre $N_0$; es claro que este último cuasi-isomorfismo es sobreyectivo.
